\begin{theorem}[\texorpdfstring{$\log\mathfrak{M}(\theta)$}{logM(theta)} 的解析性与变分公式]\label{thm:finite-part-logM-holomorphic-variation}
在小节 \ref{subsec:logM-analyticity} 的设定下，设 \(\theta\mapsto M_\theta\) 在某复邻域内全纯，并假设 \(\lambda(\theta)=\rho(M_\theta)\) 在该邻域内为简单特征值；记 \(z_\star(\theta)=\lambda(\theta)^{-1}\) 且 \(\zeta_\theta(z)=\det(I-zM_\theta)^{-1}\)。
则 \(\theta\mapsto \log\mathfrak{M}(\theta)\) 在该邻域内全纯，并且满足变分恒等式
\[
\boxed{\;
\frac{d}{d\theta}\log\mathfrak{M}(\theta)
=\frac{d}{d\theta}\log C(\theta)
\;+\;
\sum_{k\ge 2}\frac{\mu(k)}{k}\,\frac{d}{d\theta}\log \zeta_\theta\!\bigl(z_\star(\theta)^k\bigr)\; }.
\]
此外，对每个固定 \(z\)（满足 \(|z|<|z_\star(\theta)|\)）有基本微分恒等式
\[
\boxed{\;
\partial_\theta \log \zeta_\theta(z)
=
\mathrm{Tr}\!\Bigl((I-zM_\theta)^{-1}\,z\,\partial_\theta M_\theta\Bigr)\; }.
\]
\end{theorem}

\begin{proof}
由 \(\log\mathfrak{M}(\theta)=\log C(\theta)+\sum_{k\ge 2}\frac{\mu(k)}{k}\log \zeta_\theta(z_\star(\theta)^k)\) 的接口式（见 \S\ref{app:real-input-40-logm-theta}），对 \(k\ge 2\) 有 \(|z_\star(\theta)^k|<|z_\star(\theta)|\)，从而各项在邻域内全纯且尾部一致收敛，故和函数全纯并可逐项求导，得到第一式。
\par
对第二式，记 \(A_\theta(z)=I-zM_\theta\)，则 \(\log\zeta_\theta(z)=-\log\det A_\theta(z)\)。
对 \(\theta\) 求导并用 \(\partial_\theta\log\det A=\mathrm{Tr}(A^{-1}\partial_\theta A)\) 与 \(\partial_\theta A_\theta(z)=-z\,\partial_\theta M_\theta\) 即得。证毕。
\end{proof}
\leanverified{paper\_finite\_part\_logM\_holomorphic\_variation}

\begin{theorem}[\texorpdfstring{$\log\mathfrak{M}(\theta)$}{logM(theta)} 的可审计截断误差界]\label{thm:finite-part-logM-truncation-explicit-bound}
沿用定理 \ref{thm:finite-part-logM-holomorphic-variation} 的设定，并记 \(d=\dim M_\theta\)、\(\lambda=\lambda(\theta)>1\)。
定义 \(K\) 阶截断
\[
\log\mathfrak{M}_K(\theta):=\log C(\theta)+\sum_{k=2}^{K}\frac{\mu(k)}{k}\log \zeta_\theta(\lambda^{-k}).
\]
则对所有整数 \(K\ge 2\) 有显式误差上界
\[
\boxed{\;
\bigl|\log\mathfrak{M}(\theta)-\log\mathfrak{M}_K(\theta)\bigr|
\le
\frac{d}{K+1}\cdot\frac{\lambda^{-K}}{(1-\lambda^{-1})^2}\; }.
\]
\end{theorem}

\begin{proof}
写出特征值分解（按代数重数计）\(\det(I-zM_\theta)=\prod_{j=1}^{d}(1-\lambda_j z)\)，其中 \(|\lambda_j|\le\lambda\)。
当 \(k\ge 2\) 且 \(z=\lambda^{-k}\) 时，有 \(|\lambda_j z|\le \lambda^{1-k}\le\lambda^{-1}<1\)。
由不等式 \(|\log(1-w)|\le |w|/(1-|w|)\)（\(|w|<1\)）得
\[
|\log \zeta_\theta(\lambda^{-k})|
=\left|\sum_{j=1}^{d}\log(1-\lambda_j\lambda^{-k})^{-1}\right|
\le \sum_{j=1}^{d}\frac{|\lambda_j|\lambda^{-k}}{1-|\lambda_j|\lambda^{-k}}
\le d\cdot \frac{\lambda^{1-k}}{1-\lambda^{-1}}.
\]
因此
\[
\bigl|\log\mathfrak{M}(\theta)-\log\mathfrak{M}_K(\theta)\bigr|
\le
\sum_{k>K}\frac{1}{k}\,|\log \zeta_\theta(\lambda^{-k})|
\le
\frac{d}{1-\lambda^{-1}}\sum_{k>K}\frac{\lambda^{1-k}}{k}.
\]
对 \(k>K\) 有 \(1/k\le 1/(K+1)\)，故
\[
\sum_{k>K}\frac{\lambda^{1-k}}{k}
\le
\frac{1}{K+1}\sum_{k>K}\lambda^{1-k}
=
\frac{1}{K+1}\cdot\frac{\lambda^{-K}}{1-\lambda^{-1}}.
\]
合并即得所示界。证毕。
\end{proof}
\leanverified{paper\_finite\_part\_logM\_truncation\_explicit\_bound}

\begin{corollary}[截断层间的比例传递与穷竭误差]\label{cor:finite-part-logM-relative-scale-transfer}
沿用定理 \ref{thm:finite-part-logM-truncation-explicit-bound} 的记号。则对任意整数 \(L>K\ge 2\)，有
\[
\bigl|\log\mathfrak{M}_L(\theta)-\log\mathfrak{M}_K(\theta)\bigr|
\le
\frac{d}{K+1}\cdot\frac{\lambda^{-K}}{(1-\lambda^{-1})^2}.
\]
若进一步有 \(\log\mathfrak{M}(\theta)\neq 0\)，则
\[
\left|
\frac{\log\mathfrak{M}_K(\theta)}{\log\mathfrak{M}(\theta)}-1
\right|
\le
\frac{d}{(K+1)\,|\log\mathfrak{M}(\theta)|}\cdot\frac{\lambda^{-K}}{(1-\lambda^{-1})^2}.
\]
\end{corollary}

% --- Distillation writeback ---
\begin{theorem}[completed Euler denominator ledger and the commensurable--singular alternative]\label{distill:euclid_elements-proportion_and_denominator_spectrum_families-zeta-completed-denominator-ledger-alternative}
\textbf{Date-time: 2026-04-23T12:21:04+08:00. Family: Euclid--Elements / proportion and denominator-spectrum families. Dependency status: conditional audit.}
Work in the setting of Theorem \ref{thm:finite-part-logM-holomorphic-variation}, and fix one parameter value \(\theta\).  Put \(\lambda=\lambda(\theta)>1\) and write \(M=M_\theta\).  Define the completed Euler denominator ledger
\[
\widehat{\mathcal L}:=\prod_{k\ge2}\prod_{r\ge1}\mathbb C\,\mathbf e_{k,r}
\]
and the coefficient family
\[
A_{k,r}(\theta):=\frac{\mu(k)}{kr}\,\lambda^{-kr}\,\operatorname{Tr}(M^r)
\qquad(k\ge2,\ r\ge1).
\]
Then
\[
\mathfrak D_\theta:=\bigl(A_{k,r}(\theta)\bigr)_{k,r}\in\widehat{\mathcal L}
\]
is a well-defined completed ledger, and its scalar evaluation is absolutely convergent:
\[
\sum_{k\ge2}\sum_{r\ge1}|A_{k,r}(\theta)|<\infty.
\]
Moreover, if
\[
\operatorname{Supp}^{\mu}(\mathfrak D_\theta):=\bigl\{(k,r):\mu(k)\operatorname{Tr}(M^r)\ne0\bigr\},
\]
then the denominator spectrum
\[
\mathsf{Den}(\theta):=\bigl\{k\ge2:\exists r\ge1\text{ with }(k,r)\in\operatorname{Supp}^{\mu}(\mathfrak D_\theta)\bigr\}
\]
satisfies the following exact alternative.
\begin{enumerate}
  \item \textbf{Commensurable periodic case.}  If there exists \(N\in\mathbb N_{>0}\) such that \(k\mid N\) for every \(k\in\mathsf{Den}(\theta)\), then the whole non-zero Euler denominator support factors through the finite divisibility lattice of \(N\).  Equivalently, every active denominator atom is periodic with respect to the single common denominator ring \(\mathbb Z/N\mathbb Z\).
  \item \textbf{Singular finite-prime case.}  If \(\mathsf{Den}(\theta)\) is supported on a finite prime set \(P\) but is not contained in the divisor lattice of any one \(N\), then for every height bound \(H\) there is an active atom \((k,r)\in\operatorname{Supp}^{\mu}(\mathfrak D_\theta)\) such that all prime divisors of \(k\) lie in \(P\) and
  \[
  \sum_{p\in P}v_p(k)>H.
  \]
  Thus the obstruction is not the appearance of a new prime but unbounded \(P\)-adic denominator height.
  \item \textbf{Incommensurable multiprime case.}  If the union of prime divisors of elements of \(\mathsf{Den}(\theta)\) is infinite, then no finite prime-register can contain the active Euler denominator support.
\end{enumerate}
For a finite audit window \(\mathcal W\subset\{(k,r):k\ge2,r\ge1\}\), the same statements hold for the certified subledger obtained by replacing \(\operatorname{Supp}^{\mu}(\mathfrak D_\theta)\) with
\[
\operatorname{Supp}^{\mu}_{\mathcal W}(\mathfrak D_\theta)
:=\bigl\{(k,r)\in\mathcal W:\mu(k)\operatorname{Tr}(M^r)\ne0\text{ has been certified}\bigr\}.
\]
Any conclusion about the full ledger requires, in addition, zero certificates for every atom outside the certified non-zero support.  Atoms with \(\mu(k)=0\) are not denominator obstructions for this Euler ledger.
\end{theorem}
\begin{proof}
Let \(\alpha_1,\ldots,\alpha_d\) be the eigenvalues of \(M\), counted with algebraic multiplicity.  Since \(\lambda=\rho(M)\), one has \(|\alpha_j|\le\lambda\), and hence
\[
|\operatorname{Tr}(M^r)|\le\sum_{j=1}^{d}|\alpha_j|^r\le d\lambda^r.
\]
Therefore
\[
|A_{k,r}(\theta)|
\le \frac{d}{kr}\lambda^{r-kr}
=\frac{d}{kr}\lambda^{-(k-1)r}.
\]
Summing first in \(k\) and then in \(r\),
\[
\sum_{k\ge2}\sum_{r\ge1}|A_{k,r}(\theta)|
\le d\sum_{k\ge2}\frac1k\sum_{r\ge1}\frac{\lambda^{-(k-1)r}}{r}
\le d\sum_{k\ge2}\frac1k\sum_{r\ge1}\lambda^{-(k-1)r}.
\]
The inner geometric series is \(\lambda^{-(k-1)}/(1-\lambda^{-(k-1)})\), which is bounded by \(C\lambda^{-(k-1)}\) for a constant \(C=C(\lambda)\).  Thus the last series is dominated by
\[
 dC\sum_{k\ge2}\frac{\lambda^{-(k-1)}}{k}<\infty.
\]
This proves both that \(\mathfrak D_\theta\) is a legitimate element of the completed product ledger and that its scalar evaluation is absolutely convergent.  The use of \(\operatorname{Supp}^{\mu}\) is essential: if \(\mu(k)=0\), then every coefficient with that denominator is zero regardless of \(\operatorname{Tr}(M^r)\), so such an atom cannot contribute an Euler denominator obstruction.

It remains to prove the alternative.  If a common denominator \(N\) exists, then every active \(k\) belongs to the finite set \(\{d:d\mid N\}\).  Hence all active denominator comparisons are performed inside the finite divisibility lattice of \(N\), or equivalently inside the finite common denominator ring \(\mathbb Z/N\mathbb Z\).  This is precisely the commensurable periodic case.

Assume next that the active denominators use only a finite prime set \(P\), but are not contained in the divisor lattice of any fixed \(N\).  If the heights \(\sum_{p\in P}v_p(k)\) were bounded by some \(H\) on all active denominators, then only finitely many exponent vectors \((v_p(k))_{p\in P}\) could occur.  The least common multiple of the corresponding finitely many \(k\)'s would be a single integer \(N\) divisible by every active denominator, contradicting the hypothesis.  Hence for every \(H\) there is an active \(k\) supported on \(P\) with height greater than \(H\), which is exactly the singular finite-prime alternative.  Notice that the statement explicitly restricts to denominators whose prime divisors lie in \(P\); no outside-prime atom is included in this projection.

Finally, if the active prime support is infinite, then any register using only finitely many prime coordinates misses an active denominator whose prime divisor lies outside the register.  Such a support therefore cannot be represented by a finite prime-register.  The three cases are exhaustive: either the active prime support is infinite, or it is finite; in the finite case either the exponent heights are bounded, giving a common denominator by the least common multiple construction, or they are unbounded.  They are also mutually exclusive by construction.

The finite-window assertion is the same argument applied to the certified subset \(\operatorname{Supp}^{\mu}_{\mathcal W}(\mathfrak D_\theta)\).  Since uncertified atoms may later be certified non-zero, a full-ledger conclusion is justified only when the audit also supplies zero certificates outside the certified non-zero support.  This is the Omega denominator-spectrum rule: finite truncations are admissible certificates for what they explicitly contain, while global commensurability requires an all-atoms ledger audit rather than a projection illusion.
\end{proof}

% --- End distillation writeback ---

\begin{proof}
第一式中，\(\log\mathfrak{M}_L(\theta)-\log\mathfrak{M}_K(\theta)\) 正是从 \(k=K+1\) 到 \(k=L\) 的有限尾和，因此其绝对值不超过定理 \ref{thm:finite-part-logM-truncation-explicit-bound} 给出的从 \(k=K+1\) 到无穷的总尾和界。第二式则由
\[
\left|
\frac{\log\mathfrak{M}_K(\theta)}{\log\mathfrak{M}(\theta)}-1
\right|
=
\frac{\bigl|\log\mathfrak{M}_K(\theta)-\log\mathfrak{M}(\theta)\bigr|}{|\log\mathfrak{M}(\theta)|}
\]
与定理 \ref{thm:finite-part-logM-truncation-explicit-bound} 直接得到。证毕。
\end{proof}
\leanverified{paper\_finite\_part\_logM\_relative\_scale\_transfer}

\begin{remark}[穷竭法式比例读法]\label{rem:finite-part-logM-exhaustion-ratio-reading}
推论 \ref{cor:finite-part-logM-relative-scale-transfer} 表明，截断 \(\log\mathfrak{M}_K(\theta)\) 不应被看成与 \(\log\mathfrak{M}(\theta)\) 并列的另一个对象，而应被看成同一解析对象在有限层上的穷竭近似。真正稳定的是“截断层与全量对象的比例误差”而不是某个孤立的 \(K\) 层值；这正是 Euclid 式有限阶段逼近的现代版本。
\end{remark}
